\subsection{Matrix transposition over the binary prime field}\label{sec:packing:transpose}

The transpose below is an $\F_2$-linear operation on
a vector of $n$ elements of $L$.
It is generally \emph{not} an $L$-linear transformation;
this is the source of the nontrivial ring-switching
correction in the subsequent sections.

\begin{definition}[Binary coordinate matrix]\label{def:pack:matrix}
For $s=(s_v)_{v\in I}\in L^I$ define its
binary coordinate matrix
\[
 M(s)=(M_{u,v}(s))_{u,v\in I},\qquad
 M_{u,v}(s):=d_u(s_v)\in\F_2.
\]
The first index labels a binary basis coordinate;
the second labels an entry of the original field vector.
\end{definition}

\begin{lemma}[Recovering a field vector from its matrix]\label{lem:pack:matrix-reconstruct}
For every $s\in L^I$ and every $v\in I$,
\[
 s_v=\sum_{u\in I}\beta_uM_{u,v}(s).
\]
\end{lemma}
\begin{proof}
By Definition~\ref{def:pack:matrix},
$M_{u,v}(s)=d_u(s_v)$.
The displayed sum is therefore the unique
basis expansion of $s_v$, namely
$\sum_u\beta_ud_u(s_v)=s_v$.
\end{proof}

\begin{definition}[Binary coordinate transpose]\label{def:pack:transpose}
For $s\in L^I$ define $\mathcal T(s)\in L^I$ by
\[
 \boxed{\mathcal T(s)_u
      :=\sum_{v\in I}\beta_vd_u(s_v)
      =\sum_v\beta_vM_{u,v}(s).}
\]
\end{definition}

\begin{lemma}[Transpose is binary-linear]\label{lem:pack:transpose-linear}
For every $s,t\in L^I$ and
$e\in\F_2$,
\[
 \mathcal T(s+et)=\mathcal T(s)+e\mathcal T(t).
\]
\end{lemma}
\begin{proof}
For each $u$, expand
$\mathcal T(s+et)_u=
\sum_v\beta_vd_u(s_v+et_v)$.
Linearity of $d_u$ over $\F_2$ gives
$d_u(s_v)+ed_u(t_v)$.
Distribute the sum and factor out the scalar $e$.
The resulting coordinates equal those of
$\mathcal T(s)+e\mathcal T(t)$.
\end{proof}

\begin{lemma}[Coordinates of a transposed vector]\label{lem:pack:transpose-coords}
For all $u,v\in I$ and $s\in L^I$,
\[
 \boxed{d_v(\mathcal T(s)_u)=d_u(s_v).}
\]
\end{lemma}
\begin{proof}
Use Definition~\ref{def:pack:transpose}
and the coordinate property
$d_v(\beta_a)=\delta_{va}$
(Lemma~\ref{lem:binary:extract}):
\[
 d_v(\mathcal T(s)_u)
 =\sum_{a\in I}d_u(s_a)d_v(\beta_a)
 =d_u(s_v).
\]
The first equality uses $\F_2$-linearity
of $d_v$ and that $d_u(s_a)\in\F_2$.
\end{proof}

\begin{theorem}[Transposition is an involution]\label{thm:pack:transpose-involutive}
For every $s\in L^I$,
\[
 \boxed{\mathcal T(\mathcal T(s))=s.}
\]
\end{theorem}
\begin{proof}
For each $v\in I$, expand the second transpose:
\[
 \mathcal T(\mathcal T(s))_v
 =\sum_{u\in I}\beta_u
   d_v(\mathcal T(s)_u).
\]
Lemma~\ref{lem:pack:transpose-coords}
replaces the last coordinate with $d_u(s_v)$.
The sum becomes $\sum_u\beta_ud_u(s_v)$,
which equals $s_v$ by
Lemma~\ref{lem:pack:matrix-reconstruct}.
This holds for all $v$.
\end{proof}

\begin{corollary}[Transpose is a bijective binary-linear map]\label{cor:pack:transpose-iso}
$\mathcal T:L^I\to L^I$ is an $\F_2$-linear
isomorphism with inverse $\mathcal T$.
\end{corollary}
\begin{proof}
Lemma~\ref{lem:pack:transpose-linear}
gives linearity; Theorem~\ref{thm:pack:transpose-involutive}
makes the map its own two-sided inverse.
\end{proof}

\begin{lemma}[The coordinate matrices are transposes]\label{lem:pack:matrix-transpose}
For every $s\in L^I$ and all $u,v\in I$,
\[
 \boxed{M_{u,v}(\mathcal T(s))=M_{v,u}(s).}
\]
\end{lemma}
\begin{proof}
The left side is $d_u(\mathcal T(s)_v)$
by Definition~\ref{def:pack:matrix}.
Lemma~\ref{lem:pack:transpose-coords}
turns it into $d_v(s_u)$,
which is the right side.
\end{proof}

\begin{lemma}[Nonzero field elements have some nonzero coordinate]\label{lem:pack:coordinate-separates}
For every $x\in L$,
\[
 x\neq0\quad\Longleftrightarrow\quad
 \exists u\in I,\ d_u(x)=1.
\]
\end{lemma}
\begin{proof}
If every coordinate vanishes, the expansion
$x=\sum_u\beta_ud_u(x)$ is zero.
Contrapositively, $x\neq0$ implies
some coordinate is nonzero, and an element
of $\F_2$ is nonzero exactly when it is $1$.
Conversely the zero element has all coordinates
zero by $\F_2$-linearity of the maps $d_u$.
\end{proof}

\begin{proposition}[Transposition is not extension-field-linear]\label{prop:pack:not-L-linear}
Because $n=2^\ell\geq2$,
the operation $\mathcal T:L^I\to L^I$
is not $L$-linear.
\end{proposition}
\begin{proof}
Suppose toward a contradiction that $\mathcal T$
is $L$-linear.
Choose $v\in I$, and let $e_v\in L^I$
be the vector equal to $1$ at $v$ and
$0$ elsewhere. Since $1\neq0$, the preceding
lemma gives an index $u$ with $d_u(1)=1$.
For any $a\in L$, the definition of
$\mathcal T$ gives
$\mathcal T(ae_v)_u=\beta_vd_u(a)$
and $\mathcal T(e_v)_u=\beta_vd_u(1)=\beta_v$.
The assumed $L$-linearity would imply
$\beta_vd_u(a)=a\beta_v$.
Because $\beta_v\neq0$, division gives
$d_u(a)=a$ for every $a\in L$.
But $d_u(a)\in\F_2$ and $|L|=2^n>2$,
so choosing $a\in L\setminus\F_2$
is a contradiction.
\end{proof}

\begin{lemma}[Uniformity under binary transposition]\label{lem:pack:transpose-uniform}
If $S$ is uniform on $L^I$, then
$\mathcal T(S)$ is uniform on $L^I$.
\end{lemma}
\begin{proof}
By Corollary~\ref{cor:pack:transpose-iso},
$\mathcal T$ is a bijection of the finite
set $L^I$. Each output has exactly one
preimage, hence the uniform input law
is mapped to the uniform output law.
\end{proof}

\begin{remark}[Connection with field-linear proof systems]
A proof system verifying only affine equations
over $L$ cannot directly replace
$\mathcal T(m)$ by a linear combination
of hidden field entries $m_v$:
the coefficient extraction $d_u$ is only
$\F_2$-linear. Our protocol therefore
reveals masked field entries $m_v$ and
enforces the hidden correction using
affine trace coordinates, rather than
pretending that $\mathcal T$ is $L$-linear.
\end{remark}
